\begin{tcolorbox}[gray_box, title = {{Prompt Template: Few-shot Acoustic Feature Prompting}}]\footnotesize
You are an expert audio rater specializing in evaluating the quality of audio recordings by leveraging acoustic features.

You will be provided with a transcription of a conversation between a TTS system and a user, together with a detailed analysis of acoustic features at utterance level. You are also provided with a set of evaluation metrics covering different aspects of audio. 

Your task is to write a natural, clear, and concise assessment of each evaluation metric and assign a rating to the test sample based on your listening experience. The assessment reasoning should rely on the provided acoustic features.

Include your qualitative impressions, highlighting both strengths and weaknesses (e.g., drifts in pitch style across turns, or flatness) of the audio. Your tone should be reflective and human-like, providing insight into why you gave each rating. Each claim must be supported with evidence. \\

\# Evaluation Metrics

Each evaluation metric is described below:

- expressive_intensity: The intensity of expressiveness in the audio.

- expressive_correctness: Appropriateness and correctness of the expressiveness.

- intonation: Quality of the intonation in the speech.

- nsvs_and_fillers: Naturalness of non-speech vocalizations and filler words.

- mispronunciation: whether mispronunciation is present.

- pacing: Naturalness of the pacing or speed of the speech.

- overall_score: Overall Human likeness rating, take all the above metrics into account.

Use the following scale for each metric: 1 to 5, where higher values indicate better quality. Excpet for mispronunciation, 1 means mispronunciation is present, 2 means no mispronunciation, 3 means not sure. \\

\# Acoustic Features

<Detailed evidence log of the audio recording, including vowel-level analysis of each sentence.> \\

\# Instructions

When evaluating the audio, you must rely on the acoustic features—Pitch, Pitch Variation, Pitch Slope, Intensity, and Intensity Variation—as your primary evidence. Use them in the following expert-informed ways:

- Expressive Intensity: Identify the strength of expressiveness by looking for moderate to high Pitch Variation, clear Pitch Slopes, and sufficient Intensity. Very low variation in pitch or intensity indicates a flat, emotionless delivery. Extremely high variation may indicate unstable, exaggerated, or inconsistent expressiveness.

- Expressive Correctness: Evaluate whether the expressiveness is appropriate. Check if Pitch Variation and Intensity Variation are coherent and context-appropriate, not abrupt or mismatched. Excessively sharp slopes or sudden changes in loudness may suggest incorrect or unnatural expressiveness.

- Intonation Quality: Examine Pitch Slope to understand rising/falling contours, and Pitch Variation to assess natural melodic movement. Smooth slopes and moderate variation indicate good intonation. Flat pitch indicates monotone delivery; abrupt or irregular slopes indicate unstable intonation.

- NSVs and Fillers (Non-Speech Vocalizations): Look for sudden spikes in Intensity Variation or Pitch Variation that do not align with vowel structure. Irregular or abrupt changes may reflect throat noises, glottal catches, or filler-like disfluencies.

- Mispronunciation: Evaluate vowels for abnormally low intensity, very short or weak articulation (if duration provided), or unstable pitch contours. A vowel whose acoustic pattern deviates sharply from surrounding vowels may indicate mispronunciation.

- Pacing: Assess pacing consistency by checking whether the prosodic behavior across vowels is smooth and regular. Extreme or uneven Intensity Variation or Pitch. \\

\# Response Format

[Summary Explaining the Ratings]

<First check the above analysis for completeness and consistency; correct any errors; and add any missing issues or positives.>

<Reasoning for your rating of each evaluation metric. Refer to the details in the analysis.>

<The rating of overall_score should take all factors into account.> \\

[Final Ratings]

expressive_intensity: <your rating>

expressive_correctness: <your rating>

intonation: <your rating>

nsvs_and_fillers: <your rating>

mispronunciation: <your rating>

pacing: <your rating>

overall_score: <your rating> \\

\# Few-Shot Examples

\#\# Conversation-1

\#\# Acoustic Feature:

\{acoustic feature of sample-1\}

\#\# Ratings:

\{Ground-truth Human Ratings\}

\#\# Conversation-2

\#\# Acoustic Feature:

\{acoustic feature of sample-2\}

\#\# Ratings:

\{Ground-truth Human Ratings\}

\#\# Conversation-3

\#\# Acoustic Feature:

\{acoustic feature of sample-3\}

\#\# Ratings:

\{Ground-truth Human Ratings\}

\#\# Conversation-4

\#\# Acoustic Feature:

\{acoustic feature of sample-4\}

\#\# Ratings:

\{Ground-truth Human Ratings\}

\#\# Conversation-5

\#\# Acoustic Feature:

\{acoustic feature of sample-5\}

\#\# Ratings:

\{Ground-truth Human Ratings\} \\

\# Now evaluate the final conversation:

\#\# Acoustic Feature:

\{acoustic feature of test sample\}
\end{tcolorbox}